% !TeX root = ../../Book1.tex
\OptionalSection{密度与可整流性}{b1:sec:1904}

本章主线已经从参数化推出切平面、集合密度与带权可整流测度的密度。现在反过来问：如果不知道参数化，只观察小球内的质量，能否判断集合或测度可整流？这一逆问题远比顺向计算困难。本选读节先证明可整流／纯不可整流分解，再陈述作为深层输入的 Preiss 定理，最后把集合的密度判据作为推论。深层几何步骤明确交由文献，不能把它们看作前面面积公式证明的一部分。

% 实读：Mattila2023Rectifiability 出版社OA正式全书§4.3--4.4，尤其pp.23--25定理4.10--4.11及证明思路。
% Preiss1987Geometry 的Annals官方页核验作者、卷期页码DOI，但未获得其107页原文的完整证明。
% 密度判据的逆向以Preiss为深层输入；Preiss只给明确标出的proofsketch。
% 带权密度已移入19.3主线；rectifiable/pure分解为本书从有限质量上确界完整推导。

\subsection{Federer 型结构定理}
\label{b1:subsec:190403}

无论密度极限是否存在，有限 \(\mathcal H^k\) 测度的 Borel 集总能分出可被参数片捕捉的部分与完全剩余的部分。这一分解不需要 Preiss 定理。

\begin{theorem}[可整流与纯不可整流分解]\label{b1:thm:rectifiable-pure-decomposition}
设 \(E\subseteq\mathbb R^n\) 是 Borel 可测的，且 \(\mathcal H^k(E)<\infty\)。则可以写成不交并
\[
 E=R\mathbin{\dot\cup}P,
\]
其中 \(R\) 是 Borel 的且可数 \(k\) 维可整流，\(P\) 是 Borel 的且纯 \(k\) 维不可整流。这样的分解在 \(\mathcal H^k\) 零测集意义下唯一。
\end{theorem}
\begin{proof}
在所有可数 \(k\) 维可整流的 Borel 子集 \(F\subseteq E\) 中，取其测度的上确界
\[
 a=\sup\{\mathcal H^k(F):F\subseteq E\text{ 为这样的子集}\}.
\]
该数有限。取 \(F_j\) 使其测度趋于 \(a\)，并令 \(R=\bigcup_jF_j\)。可数并仍可整流，所以 \(\mathcal H^k(R)\leq a\)；又因它包含每个 \(F_j\)，有 \(\mathcal H^k(R)=a\)。置 \(P=E\setminus R\)。

任取 Lipschitz 映射 \(f:\mathbb R^k\to\mathbb R^n\)。其全空间像为 Borel 集，故 \(A=P\cap f(\mathbb R^k)\) 为 Borel 的可整流集。如果 \(\mathcal H^k(A)>0\)，则不交并 \(R\cup A\) 仍是 \(E\) 的可整流 Borel 子集，但
\[
 \mathcal H^k(R\cup A)=a+\mathcal H^k(A)>a,
\]
与上确界的定义矛盾。因此 \(P\) 纯不可整流。

若还有 \(E=R'\mathbin{\dot\cup}P'\) 满足同样条件，则 \(R\setminus R'\subseteq R\cap P'\) 为零测集，因为 \(R\) 可整流而 \(P'\) 纯不可整流。交换两组分解，得到 \(R'\setminus R\) 也零测。故 \(R\) 与 \(R'\) 的对称差零测，补集 \(P,P'\) 也只差零测集。
\end{proof}

若 \(E\) 只有 \(\sigma\) 有限测度，可先分成可数个两两不交的有限测度 Borel 块，逐块分解后分别取并。可整流部分的可数并仍可整流，纯不可整流部分与每个参数像的交则是可数个零测集之并，所以仍纯不可整流。唯一性采用同一交集论证，不需要对无穷质量作相减。

这一结构体现了 Besicovitch 的平面理论及 Federer 的高维推广中最基本的区分：参数化所能识别的部分可以集中保留，剩下的部分与每张参数片都只零测相交。完整的 Federer 型结构理论还研究投影、近似切平面与纯不可整流部分之间的关系。那些投影判据需要另外指定“几乎每个方向”所用的测度，不应从上面的上确界分解自行推出。可进一步参阅 Federer 的《Geometric Measure Theory》\cite{Federer1996Geometric}及 Mattila 的《Rectifiability: A Survey》\cite[Section 4]{Mattila2023Rectifiability}。

\subsection{Preiss 定理概览}
\label{b1:subsec:190404}

现在陈述从密度重建可整流结构所需的深层结果，即\term[preiss-dingli]{Preiss 定理}[Preiss's theorem]。与集合测度不同，一般 Radon 测度的密度不必等于 \(1\)，因此正确假设是密度存在、正且有限。

\begin{theorem}[Preiss]\label{b1:thm:preiss-rectifiability}
设 \(\mu\) 为 \(\mathbb R^n\) 上的正 Radon 测度，\(1\leq k\leq n\) 为整数。如果
\[
 0<\lim_{r\downarrow0}\frac{\mu(B(x,r))}{\omega_kr^k}<\infty
 \quad\text{对 }\mu\text{-几乎每个 }x,
\]
则 \(\mu\) 是本书意义下的 \(k\) 维可整流测度。反过来，每个这样的可整流测度都满足上述密度条件。
\end{theorem}

定理的顺向几何结论来源于 Preiss 的《Geometry of Measures in \(\mathbf R^n\): Distribution, Rectifiability, and Densities》\cite{Preiss1987Geometry}。本书所述版本可与 Mattila 的《Rectifiability: A Survey》\cite[Theorem 4.11 and the following discussion]{Mattila2023Rectifiability}核对；其中从有限上密度得到 \(\mu\ll\mathcal H^k\)，所以也满足本书较强的测度约定。反向已经由\cref{b1:prop:rectifiable-weight-density}完整证明。下面只解释正向证明的组织，不尝试用几段话替代原论文的长证明。

\begin{proofsketch}
先看缩放。对 \(x\in\mathbb R^n\)、\(r>0\)，令 \(T_{x,r}(z)=(z-x)/r\)。把 \(\mu\) 推前并重新调整质量后，考察
\[
 c_j(T_{x,r_j})_\#\mu,
 \quad c_j>0,\quad r_j\downarrow0,
\]
在 \(C_c(\mathbb R^n)\) 测试下的非零 Radon 极限。这样的极限称为 \(\mu\) 在 \(x\) 处的\term[qiecedu]{切测度}[tangent measure]。它是一个测度，而不是一张线性平面；其收敛是\chapref{b1:ch:12}的淡收敛。

正的有限密度给缩放测度提供局部质量控制。以 \(c_j=r_j^{-k}\) 为例，每个固定半径球的质量保持有界，原点附近又不会全部消失。这使提取非零局部极限成为可能。仅有这一紧性还不够，重要的是把原测度在几乎处处的密度信息转移到这些极限的全部支撑点，得到统一的球质量关系。

接下来需要研究满足该关系的测度的结构。这里有一个不可省略的困难：当维数较高时，具有同一幂次球质量的测度不一定集中在一张平面上。因此不能在得到这种极限后直接宣布“切测度是平面的”。Preiss 的证明还使用切测度族的结构和进一步缩放，从中取得足够的平坦性，再回到原测度，构造可数个捕捉其质量的参数片。

本概要省略了局部极限的系统构造、统一球质量的传递、非平坦极限的排除机制以及最后的可整流性恢复；这些正是定理的深层部分，而非例行估计。Mattila 的《Rectifiability: A Survey》\cite[Sections 4.3 Tangent Measures and 4.4 Densities]{Mattila2023Rectifiability}给出这些步骤之间的关系和进一步文献，完整论证须读 Preiss 原文。这里没有把任意一个切测度的存在误当成已经得到可整流性。
\end{proofsketch}

切平面与切测度的区别也可以在最熟悉的例子上看出。若 \(\mu=\theta\mathcal H^k|_V\)，其中 \(V\) 为平面、\(\theta>0\) 为常数，则在 \(x\in V\) 处，\(r^{-k}(T_{x,r})_\#\mu=\theta\mathcal H^k|_{V-x}\)。平面给出方向，测度另外保留质量常数；更一般的缩放极限还可能保留重数和非平坦形状。

Preiss 定理的结论是测度意义下可数参数化，并不把承载集变成处处光滑的流形，不保证全局一张图，也不强迫密度为整数。它说明，小球质量在全部小尺度上稳定到一个正有限幂律时，欧氏几何中会出现可整流结构。

\subsection{可整流性的密度判据}
\label{b1:subsec:190402}

顺向结论只用了微分、面积与测度微分。逆向则需要从小球质量中重建方向；有了上一小节的 Preiss 定理，现在可以把这一步作为正式推论写出。

\begin{theorem}[可整流性的密度判据]\label{b1:thm:rectifiability-density-criterion}
设 \(E\subseteq\mathbb R^n\) 是 Borel 可测的，\(1\leq k\leq n\)，且 \(\mathcal H^k(E)<\infty\)。下列条件等价：
\begin{equivalences}
\item\label{b1:item:density-criterion-rectifiable}\(E\) 可数 \(k\) 维可整流；
\item\label{b1:item:density-criterion-one}\(\Theta^k(\mathcal H^k|_E,x)=1\) 对 \(\mathcal H^k\)-几乎每个 \(x\in E\) 成立；
\item\label{b1:item:density-criterion-exists}\(\Theta^k(\mathcal H^k|_E,x)\) 对 \(\mathcal H^k\)-几乎每个 \(x\in E\) 存在。
\end{equivalences}
\end{theorem}
\begin{proof}[由 Preiss 定理推出]
第一项推出第二项由\cref{b1:thm:rectifiable-density-tangent}给出，第二项当然推出第三项。若第三项成立，由\chapref{b1:ch:15}的 Hausdorff 上密度估计，
\[
 2^{-k}\leq\Theta^{*k}(\mathcal H^k|_E,x)\leq1
\]
在 \(E\) 中几乎处处成立。因此一旦密度极限存在，它必正且有限。将\cref{b1:thm:preiss-rectifiability}应用于有限 Radon 测度 \(\mu=\mathcal H^k|_E\)，得到 \(\mu\) 集中在可数个 Lipschitz 像上。因而这些像在 \(E\) 中遗漏的 \(\mathcal H^k\) 测度为零，这正是第一项。
\end{proof}

这里已经完整写出各条件之间的推导，但逆向所调用的 Preiss 定理并未在本书内证明。Mattila 的《Rectifiability: A Survey》\cite[Theorems 4.10--4.11]{Mattila2023Rectifiability}将这两层结果相邻表述，便于核对有限测度、几乎处处及欧氏空间等假设。前 3 节关于参数化和切向面积的结论都不反向依赖这条深层判据。

单有正的有限上密度不够：\chapref{b1:ch:15}已经证明这一性质对任意有限 \(\mathcal H^k\) 测度的 Borel 集几乎处处成立，而四角 Cantor 集提供了纯不可整流的例子。因此新的要求不是某个尺度上存在面积，而是所有趋零尺度上的归一化质量具有同一个极限。它把无穷多个尺度联系起来，不能用某一列半径上的收敛替代。至此，本编从覆盖与微分走到了低维集合的局部形状；下一编将用弱导数，并在 Sobolev 空间和有界变差函数空间中研究函数，届时这些几何工具会重新出现在迹、边界和有限周长问题中。
